\chapter*{Abstract}
\addcontentsline{toc}{chapter}{Abstract}

Customer reviews published on digital platforms provide an important source of insight into brand perception, service quality, and customer satisfaction. For app-based service companies such as Zomato, Google Play Store reviews reflect user experiences related to delivery, refunds and cancellations, customer support, pricing, trust, food quality, and app usability. However, because such reviews are large in volume and unstructured in nature, extracting business insights through manual analysis is difficult and inefficient.

This study analyzes customer sentiment toward Zomato using Google Play Store reviews and examines the managerial implications of the findings for brand strategy and customer experience. The study uses a cleaned and deduplicated dataset of Zomato app reviews collected from the Google Play Store. Sentiment labels are derived using a rating-based proxy in which low ratings are mapped to negative sentiment, medium ratings to neutral sentiment, and high ratings to positive sentiment. In addition, thematic classification is carried out through an ordered rule-based keyword approach to identify the major issues shaping customer opinion.

The analysis is designed to classify reviews into positive, negative, and neutral categories, identify dominant themes affecting customer perception, and interpret how these patterns can support managerial decision-making. The study contributes to practice by showing how app review analytics can be used to convert large-scale customer feedback into actionable insights for improving service performance, brand perception, and customer engagement. A limited competitor comparison is also included at the discussion stage to provide broader market context.

\textbf{Keywords:} Zomato, sentiment analysis, Google Play Store reviews, natural language processing, customer experience, brand strategy
